% !TeX root = ../../20260524.tex
\section{Applications}

\subsection{常見複雜度}

\begin{frame}{常見複雜度}
    學會計算複雜度後
    
    我們可以透過觀察題目給的範圍去判斷跑的跑不動\pause

    對於現代機器來講，通常 1s 可以跑 $5 \times 10^8$ 左右
\end{frame}

\begin{frame}{常見複雜度}
    \begin{itemize}
        \item $N \le 10^{18}$ : $O(\log{N})$\pause
        \item $N \le 10^8$ : $O(N)$\pause
        \item $N \le 10^6$ : $O(N \log{N})$\pause
        \item $N \le 10^5$ : $O(N \sqrt{N})$\pause
        \item $N \le 10^3$ : $O(N^2)$\pause
        \item $N \le 500$ : $O(N^3)$\pause
        \item $N \le 100$ : $O(N^4)$\pause
        \item $N \le 20$ : $O(2^N)$\pause
        \item $N \le 8$ : $O(N!)$
    \end{itemize}
\end{frame}

\begin{frame}{跑的跑不動?}
    要背上面的表嗎?\pause

    忘記的話，我比賽現場要怎麼知道會過?\pause

    把數字代進去。
\end{frame}

\subsection{級數複雜度}

\begin{frame}{級數複雜度}
    當然，世界不是那麼美好的

    有時候有些時間複雜度有點難分析\pause

    $O(1) + O(1) + O(2) + \cdots + O(N)$?\pause

    等差級數

    $\displaystyle O(\frac{(1 + N) \times N}{2}) = O(N^2)$
\end{frame}

\begin{frame}{級數複雜度}
    $O(1) + O(2) + O(4) + \cdots + O(2^N)$?\pause

    等比級數

    $\displaystyle O(\frac{2^{N+1}-1}{1}) = O(2^N)$
\end{frame}

\begin{frame}{級數複雜度}
    $O(N) + O(\frac{1}{2}N) + O(\frac{1}{3}N) + \cdots + O(\frac{1}{N}N)$?\pause

    $O(N \log{N})$\pause

    $O(N) + O(\frac{1}{2}N) + O(\frac{1}{4}N + \cdots + O(\frac{1}{2^k}N)), 2^k \le N$?\pause

    $O(N)$
\end{frame}

\subsection{常數}

\begin{frame}{常數}
    時間複雜度常常忽略常數

    是因為常數往往影響很小，在變數的快增長下顯得不重要\pause

    但是有時候，常數也可能導致 TLE 

    像是 $O(N\log{N})$ 的複雜度可以過

    但是你外面帶個常數 $67$，就倒了。
\end{frame}